\chapter{Security Hardening Checklist}
\label{ch:security-checklist}

Chapters~\ref{ch:threats} and~\ref{ch:access-control} in one page.

Start with the architectural section. Everything after it is defence in depth,
and defence in depth on top of a bad architecture is decoration.

\section{Architecture: the trifecta}

\begin{itemize}[leftmargin=1.6em, itemsep=3pt]
\item[$\square$] \textbf{List your agent contexts.} Each loop is one.
\item[$\square$] \textbf{For each, mark the three ingredients:} access to private
data, exposure to untrusted content, and an exfiltration channel.
\item[$\square$] \textbf{Any context with all three needs a named control.} Not
``we are careful''. A specific, testable one.
\item[$\square$] \textbf{Split privileges} so the reading context and the sending
context are not the same context.
\item[$\square$] \textbf{Write a test for each control}, or it will be
refactored away. (\chapref{threats})
\end{itemize}

\section{Host}

\begin{itemize}[leftmargin=1.6em, itemsep=3pt]
\item[$\square$] Consent prompts show \textbf{arguments}, not just tool names.
(\chapref{building-hosts})
\item[$\square$] Standing approvals are \textbf{scoped to a tool}, never
``allow everything''.
\item[$\square$] Related prompts are batched, so users do not learn to click.
\item[$\square$] Tool definitions are \textbf{pinned after approval}; a changed
hash re-prompts. (\chapref{threats})
\item[$\square$] Tool annotations from untrusted servers are \textbf{not}
believed. (\chapref{tools})
\item[$\square$] Auto-approved read-only tools are reviewed for what private data
they can read. (\chapref{threats})
\item[$\square$] Tool calls come from \textbf{structured output}, never from
parsing text.
\item[$\square$] Arbitrary \texttt{https://} resource URIs are \textbf{not}
auto-fetched. (\chapref{threats})
\item[$\square$] Step, token, and cost budgets are set to numbers.
(\chapref{agent-loop})
\item[$\square$] An unhealthy server drops out of the catalogue without taking
the host with it. (\chapref{building-hosts})
\item[$\square$] The default input provider \textbf{declines}.
\end{itemize}

\section{Server}

\begin{itemize}[leftmargin=1.6em, itemsep=3pt]
\item[$\square$] Every tool input validated against its schema at the edge.
(\chapref{building-servers})
\item[$\square$] Remote \texttt{\$ref} \textbf{never} dereferenced. Schemas that
fail to resolve are rejected, not treated as permissive. (\chapref{tools})
\item[$\square$] Schema depth and subschema count bounded. (\chapref{tools})
\item[$\square$] Tool visibility filtered by \textbf{authorization}, never by
connection. (\chapref{building-servers})
\item[$\square$] Handlers check authorization too. Filtering the catalogue is a
menu, not a lock. (\chapref{access-control})
\item[$\square$] Forbidden and nonexistent return the \textbf{same error}, or
your errors enumerate your data. (\chapref{access-control})
\item[$\square$] Handles are opaque, high-entropy, expiring, and authorized on
every use. (\chapref{building-servers})
\item[$\square$] Resource reads authorized per URI; file paths sanitised against
traversal. (\chapref{resources})
\item[$\square$] Mutating tools take an idempotency key, because retries are
routine. (\chapref{transport})
\item[$\square$] Tool outputs sanitised, and rate limits applied.
\item[$\square$] Extensions declared only when implemented.
(\chapref{building-servers})
\end{itemize}

\section{MRTR and \texttt{requestState}}

\begin{itemize}[leftmargin=1.6em, itemsep=3pt]
\item[$\square$] \textbf{Integrity protected} with a MAC, verified in constant
time.
\item[$\square$] \textbf{Bound to the principal}, blocking cross-user replay.
\item[$\square$] \textbf{Bound to the request}, blocking cross-request replay.
\item[$\square$] \textbf{Expiring}, with a short TTL.
\item[$\square$] Signed with a \textbf{fleet-wide} secret, not per process.
\item[$\square$] Rotation window longer than the state TTL.
(\chapref{operations})
\item[$\square$] \textbf{No secrets inside.} HMAC gives integrity, not
confidentiality.
\item[$\square$] Single-use enforced \textbf{server-side} if the operation
requires it. Binding and expiry do not provide it.
\item[$\square$] Never send an input request the client did not declare support
for. (\chapref{mrtr})
\item[$\square$] Form mode \textbf{never} for passwords, keys, tokens, or payment
credentials. URL mode instead. (\chapref{mrtr})
\item[$\square$] URL mode verifies that the user who \emph{opens} the URL is the
user it was minted for. (\chapref{mrtr})
\end{itemize}

\section{Transport}

\begin{itemize}[leftmargin=1.6em, itemsep=3pt]
\item[$\square$] \texttt{Origin} validated; \texttt{403} when disallowed.
(\chapref{transport})
\item[$\square$] Local servers bind to \texttt{127.0.0.1}, not
\texttt{0.0.0.0}.
\item[$\square$] Header and body agreement enforced, \texttt{400} with
\texttt{-32020} on mismatch. This is the request-smuggling control.
\item[$\square$] Base64 sentinel encoding applied, including to values that
merely resemble the sentinel.
\item[$\square$] Sensitive parameters \textbf{not} marked \texttt{x-mcp-header};
headers are visible to intermediaries. (\chapref{tools})
\item[$\square$] TLS everywhere that is not a local pipe.
\item[$\square$] stdio servers get an explicit environment allowlist.
(\chapref{operations})
\end{itemize}

\section{Authorization}

\begin{itemize}[leftmargin=1.6em, itemsep=3pt]
\item[$\square$] Protected resource metadata published at
\texttt{/.well-known/\allowbreak oauth-protected-\allowbreak resource}.
(\chapref{authorization})
\item[$\square$] Audience binding (RFC 8707) enforced. Tokens issued for someone
else are rejected.
\item[$\square$] \textbf{No token passthrough.} Never forward the client's token
to a third party.
\item[$\square$] \texttt{iss} validated before redeeming a code, by simple string
comparison, with no normalisation, including on error responses.
\item[$\square$] Scope challenges include everything the operation needs, in one
challenge.
\item[$\square$] Clients request the \textbf{union} on step-up, not just the new
scope.
\item[$\square$] Re-auth retries are bounded.
\item[$\square$] Client credentials keyed by issuer, never reused across
authorization servers.
\end{itemize}

\section{MCP Apps}

\begin{itemize}[leftmargin=1.6em, itemsep=3pt]
\item[$\square$] Templates reviewed before rendering; predeclaration exists to
make this possible. (\chapref{apps})
\item[$\square$] Declared CSP origins read and allowlisted. A \texttt{connect-src}
to an unknown host is an exfiltration request in writing.
\item[$\square$] Network denied by default.
\item[$\square$] \texttt{postMessage} origin and payload shape validated on every
message.
\item[$\square$] \textbf{UI-initiated tool calls traverse the normal consent and
audit path.} No exceptions.
\item[$\square$] Tools degrade to useful text for hosts that render nothing.
\end{itemize}

\section{Multi-agent}

\begin{itemize}[leftmargin=1.6em, itemsep=3pt]
\item[$\square$] Delegation depth checked \textbf{before} doing work.
(\chapref{multi-agent})
\item[$\square$] Budgets are \textbf{remaining and decrementing}, not per-node
allowances, or fan-out multiplies them.
\item[$\square$] An agent with no budget assumes the \textbf{smallest}, never
unlimited.
\item[$\square$] Children act with \textbf{their own} scoped credentials.
\item[$\square$] Circuit breakers, with a half-open probe.
\item[$\square$] \texttt{traceparent} propagated through every hop.
\end{itemize}

\section{Audit}

\begin{itemize}[leftmargin=1.6em, itemsep=3pt]
\item[$\square$] Append-only, and in a different store from application data.
(\chapref{access-control})
\item[$\square$] Written \textbf{before} the action, so timeouts are recorded.
\item[$\square$] Principal from the \textbf{token}, never from an argument.
\item[$\square$] \texttt{traceparent} recorded, so one incident is one query.
\item[$\square$] Approver recorded by \textbf{name} for human-gated operations.
\item[$\square$] \textbf{Inputs recorded, not only actions}: a digest, source, and
length of what entered the model's context. Without this an injection is
invisible in the record. (\chapref{threats})
\item[$\square$] Hashes rather than content, so the log does not become a second
copy of the data.
\item[$\square$] Retention longer than your realistic detection time.
\item[$\square$] \texttt{clientInfo} logged but \textbf{never} used for
decisions.
\end{itemize}

\section{The twelve questions}

Run these against a live deployment. Each has a yes-or-no answer, and each
answer is a finding or a relief.

\begin{enumerate}[leftmargin=1.6em, itemsep=2pt]
\item Does any agent context hold private data, untrusted content, and a way out?
\item Can a tool result reach the model without being marked as data?
\item Which private data can your auto-approved read-only tools reach?
\item Can the host fetch an arbitrary \texttt{https://} resource URI?
\item Are tool definitions pinned after approval?
\item Is \texttt{requestState} integrity-protected, principal-bound, and
request-bound? All three?
\item Can a compromised server reach the internet?
\item Does the audit log record what entered the model's context?
\item Are tool calls ever extracted by parsing text?
\item Do MCP App tool calls traverse the consent path?
\item Which servers' annotations do you trust, and why?
\item What happens when a server returns 40\,MB?
\end{enumerate}

Question three finds more real problems than the rest combined, because
auto-approving read-only tools is so obviously reasonable and ``read-only'' does
not mean ``harmless''.
